%!TEX root = ../main.tex
\chapter{Background}
\label{cap:background}

\section{Reinforcement Learning}
\label{sec:background-rl}


MEMO PER ME : DA SPOSTARE!!!!!!!!!!!!!!!!!!!!!!!!!!!!!!!!!!!!!!!!!!!!!!!!!!!!!!!!!!!!!!!!!!!!!!!!!!!!!!!!!!!!!!!!!!!!!!!!-----
--------------------------------------------------------------------------------------------------------------------------------

LSTM con pesi $W_g \in
\mathbb{R}^{2d_1 \times d_1}$ e bias $b_g \in \mathbb{R}^{d_1}$ condivisi da
ogni nodo della rete, per ciascuno dei quattro gate $g \in \{\text{forget,
input, output, cell}\}$:
\begin{equation}
  \begin{bmatrix} f_i^t \\ \iota_i^t \\ o_i^t \\ g_i^t \end{bmatrix} =
  \begin{bmatrix} \sigma \\ \sigma \\ \sigma \\ \tanh \end{bmatrix}
  \left(
    W_{\{f,\iota,o,g\}}^{\top} [\,e_i^t \,\|\, h_i^{t-1}\,] + b_{\{f,\iota,o,g\}}
  \right), \qquad
  c_i^t = f_i^t \odot c_i^{t-1} + \iota_i^t \odot g_i^t, \qquad
  x_i^t = h_i^t = o_i^t \odot \tanh(c_i^t)
  \label{eq:stgat-lstm}
\end{equation}

-------------------------------------------------------------------------------------------
-------------------------------------------------------------------------------------------





Un problema di decisione sequenziale in ambiente stocastico è formalizzato 
come processo decisionale di Markov (MDP) \citep{suttonbarto2018}, definito dalla tupla
\[\mathcal{M} = (\mathcal{S}, \mathcal{A}, P, R, \gamma)\] dove $\mathcal{S}$ è
lo spazio degli stati, $\mathcal{A}$ lo spazio delle azioni, $P(s' \mid s,a)$
la probabilità di transizione markoviana,
$R(s,a)$ la ricompensa attesa e $\gamma \in [0,1)$ il fattore di sconto, che
pondera l'importanza delle ricompense future rispetto a quelle immediate. 
Un agente segue una politica
$\pi(a \mid s)$, una distribuzione di probabilità sulle azioni condizionata
allo stato corrente; il suo obiettivo è massimizzare il ritorno atteso
\[G_t = \sum_{k=0}^{\infty} \gamma^k r_{t+k}\] ovvero la somma scontata delle
ricompense future a partire dall'istante $t$. La funzione valore-stato
$V^\pi(s) = \mathbb{E}_\pi[G_t \mid s_t = s]$ e la funzione valore-azione
$Q^\pi(s,a) = \mathbb{E}_\pi[G_t \mid s_t=s, a_t=a]$ misurano rispettivamente
quanto ritorno atteso si ottiene seguendo $\pi$ da uno stato, o da una
coppia stato-azione. Queste soddisfano l'equazione di Bellman, che esprime il
valore di uno stato in funzione del valore atteso degli stati successivi:
\begin{equation}
  V^\pi(s) = \sum_a \pi(a\mid s) \sum_{s'} P(s'\mid s,a) \big[R(s,a) + \gamma V^\pi(s')\big]
\end{equation}
Per una politica ottima $\pi^*$ che massimizza $V^\pi(s)$ per ogni
stato simultaneamente, vale l'analoga equazione di Bellman di ottimalità:
\begin{equation}
  Q^*(s,a) = \mathbb{E}_{s' \sim P(\cdot \mid s,a)}\left[
    R(s,a) + \gamma \max_{a'} Q^*(s', a')
  \right]
  \label{eq:bellman-optimal}
\end{equation}
Il Q-learning tabellare approssima $Q^*$ calcolando iterativamente 
l'Equazione~\ref{eq:bellman-optimal}, ma richiede una tabella con tutte le coppie 
$(s,a)$, il che risulta inapplicabile per spazi di stato $\mathcal{S}$ 
continui o ad alta dimensionalità, come nel caso del controllo semaforico.

\paragraph{Deep Q-Network (DQN).} \citet{mnih2015dqn} superano tale limite sostituendo la tabella con
una rete neurale $Q(s,a;\theta)$. Inoltre, per mitigare l'instabilità dell'addestramento dovuta alla correlazione temporale 
dei campioni, introducono il \emph{replay buffer}, ovvero una memoria da cui vengono campionati mini-batch 
di transizioni $(s,a,r,s')$ in modo casuale, ripristinando l'assunzione di indipendenza 
necessaria per la discesa del gradiente stocastico.

\paragraph{Esplorazione $\varepsilon$-greedy.} Un agente che seguisse sempre
e solo l'azione a valore $Q$ stimato più alto non raccoglierebbe mai
esperienza sulle conseguenze di azioni alternative, rischiando di convergere
su una politica sub-ottima per via di una mancata esplorazione. 
La strategia più comune per bilanciare sfruttamento (\emph{exploitation},
scegliere l'azione ritenuta migliore) ed esplorazione
(\emph{exploration}, raccogliere informazione su azioni meno note) è
$\varepsilon$-greedy: a ogni passo di decisione, con probabilità
$\varepsilon$ l'agente sceglie un'azione uniformemente a caso, e con
probabilità $1-\varepsilon$ sceglie $\operatorname*{argmax}_a Q(s,a;\theta)$.
Il valore di $\varepsilon$ tipicamente decresce nel corso dell'addestramento, da
un valore iniziale vicino a 1, dunque da un'esplorazione quasi totale, 
verso un valore finale piccolo ma non nullo, dunque conducendo ad un comportamento quasi del tutto deterministico.

\paragraph{Estensioni usate in questo lavoro.} Il DQN di base ammette diverse
estensioni indipendenti, ciascuna delle quali corregge una specifica fonte di
instabilità o inefficienza; la procedura di training descritta nel
Capitolo~\ref{cap:tsc-modelli} le adotta tutte, motivandole caso per caso
rispetto all'articolo di riferimento:
\begin{itemize}
  \item Il \textbf{Double DQN} \citep{vanhasselt2016double} disaccoppia la
    selezione dell'azione dalla sua valutazione, usando la rete online per
    scegliere l'azione migliore e una \emph{target network} $Q(\cdot;\theta^-)$ solo per valutarla.
    Tale rete viene aggiornata periodicamente con i pesi della rete online e viene 
    usata per calcolare il target $y = r + \gamma \max_{a'} Q(s',a';\theta^-)$, in particolare  
    $y = r + \gamma\, Q\!\left(s', \operatorname*{argmax}_{a'} Q(s',a';\theta);\ \theta^-\right)$.
    Questo meccanismo evita il problema del \emph{moving target}, il quale si verifica invece ottimizzando questa stima rispetto 
    a parametri in continua variazione, e riduce il bias sistematico di sovrastima del DQN standard.
  \item Il \textbf{Prioritized Experience Replay} (PER,
    \citealp{schaul2016per}) sostituisce il campionamento uniforme dal
    replay buffer con un campionamento proporzionale all'errore di
    differenza temporale (\emph{TD error}) $\delta_i$, ovvero la discrepanza tra
    il ritorno futuro stimato e quello ottenuto. Assegnando a ciascuna transizione una priorità
    $p_i = |\delta_i| + \varepsilon$, la probabilità di campionamento della stessa è
    $P(i) = p_i^\alpha / \sum_k p_k^\alpha$, con $\alpha \in [0,1]$ iperparametro che ne controlla la dipendenza dall'errore. Per correggere il bias
    introdotto da un campionamento non uniforme, la loss viene pesata tramite coefficienti di \emph{importance sampling} 
    $w_i = \left(N \cdot P(i)\right)^{-\beta}$, dove $N$ è la dimensione del buffer e $\beta \in [0,1]$ iperparametro dinamico che ne controlla il grado di correzione.
  \item La \textbf{Huber loss} (o \emph{smooth L1}) sostituisce l'errore
    quadratico medio (MSE) con una funzione quadratica vicino allo zero e lineare
    lontano da esso, risultando meno sensibile agli outlier prodotti da transizioni con
    errore di stima molto grande:
    \begin{equation}
      L_\delta(y, \hat y) =
      \begin{cases}
        \tfrac{1}{2}(y-\hat y)^2 & \text{se } |y-\hat y| \le \delta \\
        \delta\left(|y-\hat y| - \tfrac{1}{2}\delta\right) & \text{altrimenti}
      \end{cases}
    \end{equation}
    Tipicamente $\delta = 1$.
  \item Il \textbf{soft update} (\emph{Polyak averaging}) sostituisce l'aggiornamento 
    periodico discreto della \emph{target network} con una media mobile esponenziale: 
    \[\theta^- \leftarrow \tau\,\theta + (1-\tau)\,\theta^-\] con $\tau \ll 1$, 
    stabilizzando le variazioni del target a ogni passo di ottimizzazione.
  \item Quando la rete $Q$ include un modulo ricorrente (ad esempio una LSTM come nel nostro caso, vedi Capitolo~\ref{cap:tsc-modelli}), 
    l'assunzione di Markov sulle singole transizioni campionate dal \emph{replay buffer} viene invalidata.
    Infatti, lo stato ricorrente $(h,c)$ usato durante il training viene inizializzato a
    zero, introducendo un disallineamento rispetto allo stato generato sequenzialmente durante la raccolta dati: un problema noto come \emph{hidden state
    staleness}. Seguendo \citet{kapturowski2019r2d2}, campioniamo intere
    sequenze di lunghezza $L$ e dedichiamo i primi passi della sequenza (\emph{burn-in}) esclusivamente a
    ricostruire uno stato ricorrente plausibile, senza calcolare gradiente,
    prima di eseguire la \emph{backpropagation through time} (BPTT) sui
    passi restanti.
\end{itemize}

\section{Graph Neural Network}
\label{sec:background-gnn}

Numerosi domini applicativi generano dati privi di una struttura regolare (come le immagini) o sequenziale (come il testo). 
Questi dati si organizzano spesso in topologie a grafo dove le relazioni spaziali veicolano informazione cruciale.
Un esempio naturale di grafo è costituito da una rete stradale: le intersezioni rappresentano i nodi, mentre le strade che le collegano sono gli archi. 
Lo stato di un'intersezione (ad esempio, il numero di veicoli in attesa o la fase semaforica attiva) è influenzato da ciò che accade nelle intersezioni adiacenti.
Le \emph{Graph Neural Network} (GNN) estendono i paradigmi di convoluzione 
e attenzione a questi domini a grafo, implementando meccanismi di \emph{message passing} tra nodi adiacenti.

\paragraph{Graph Convolutional Network (GCN).} \citet{kipf2017gcn} propongono la
formulazione più basilare di propagazione su grafo tra quelle qui
discusse: la rappresentazione di ciascun nodo si aggiorna mediando quelle
dei vicini con un peso fisso, determinato unicamente dalla topologia locale,
non dal contenuto delle rappresentazioni stesse. Indicando con
$\hat A = A + I$ la matrice di adiacenza a cui è stato aggiunto un
self-loop per ogni nodo (cosicché un nodo contribuisca anche al proprio
aggiornamento) e con $\hat D$ la matrice diagonale dei gradi corrispondente
($\hat D_{ii} = \sum_u \hat A_{iu}$), la regola di propagazione di un
singolo layer si scrive, in forma matriciale:
\begin{equation}
  H^{(l+1)} = \sigma\!\left(\hat D^{-1/2} \hat A\, \hat D^{-1/2} H^{(l)} W^{(l)}\right)
  \label{eq:gcn-propagation}
\end{equation}
dove $H^{(l)} \in \mathbb{R}^{I \times d_l}$ raccoglie le rappresentazioni
di tutti gli $I$ nodi del grafo al layer $l$, $W^{(l)}$ è la matrice di
pesi condivisa da ogni nodo e $\sigma$ una non linearità puntuale. Per un
singolo nodo $i$, la normalizzazione simmetrica
$\hat D^{-1/2}\hat A\hat D^{-1/2}$ equivale a pesare il contributo di ogni
vicino $u$ con $1/\sqrt{\hat d_i \hat d_u}$:
\begin{equation}
  h_i' = \sigma\!\left(\sum_{u \in \mathcal{N}(i) \cup \{i\}}
    \frac{1}{\sqrt{\hat d_i \hat d_u}}\, W h_u\right)
  \label{eq:gcn-aggregate}
\end{equation}
un peso identico per ogni arco del vicinato di $i$, funzione solo del grado
dei due nodi coinvolti e mai del loro contenuto. È, per costruzione, il
termine di paragone più naturale rispetto a cui misurare cosa aggiunge un
meccanismo di attenzione appreso come quello discusso di seguito.


\paragraph{Graph Attention Network (GAT).} \citet{velickovic2018gat} sostituiscono
il peso fisso della GCN con un coefficiente di attenzione appreso,
specifico per ciascuna coppia di nodi collegati: il contributo di un
vicino non dipende più soltanto dalla topologia, ma anche dal suo
contenuto. Dato un nodo $i$ con vicinato $\mathcal{N}(i)$, embedding di
input $h_i$, matrice di proiezione condivisa $W$ e vettore di attenzione
$\vec a$:
\begin{equation}
  e_{iu} = \mathrm{LeakyReLU}\!\left(\vec a^\top [W h_i \,\|\, W h_u]\right),
  \qquad
  \alpha_{iu} = \frac{\exp(e_{iu})}{\sum_{u' \in \mathcal{N}(i)} \exp(e_{iu'})}
  \label{eq:gat-attention}
\end{equation}
\begin{equation}
  h_i' = \sigma\!\left(\sum_{u \in \mathcal{N}(i)} \alpha_{iu}\, W h_u\right)
  \label{eq:gat-aggregate}
\end{equation}
dove $\|$ denota la concatenazione. Con $K$ teste di attenzione
indipendenti, gli output si concatenano nei layer intermedi e si mediano
nell'ultimo, secondo lo schema multi-head introdotto per l'attenzione da
\citet{vaswani2017attention}. A differenza della GCN, il coefficiente
$\alpha_{iu}$ consente a un nodo di ponderare diversamente i propri vicini
in base al loro contenuto, non solo alla loro posizione nel grafo: una
proprietà rilevante per un'intersezione i cui vicini abbiano importanza
marcatamente diversa, come nel caso di un'arteria principale contrapposta
a una strada secondaria. La funzione di compatibilità additiva appena
descritta, per quanto la più nota, non è l'unica possibile: il
Capitolo~\ref{cap:tsc-modelli} (Meta-GAT) ne impiega una variante
strutturalmente diversa, con Query, Key e Value esplicite e distinte al
posto della proiezione $W$ condivisa, e un prodotto scalare al posto della
concatenazione.

\paragraph{Reti spatio-temporali su grafo.} Una rete stradale non è
soltanto un grafo: è un grafo il cui stato (occupazione delle corsie, fase
attiva) evolve nel tempo, e le due dimensioni risultano intrecciate, poiché
il valore osservato a un nodo in un dato istante è correlato sia con i
valori passati allo stesso nodo sia con quelli ai nodi limitrofi. La
famiglia di modelli che combina un layer di propagazione sul grafo (GCN o
GAT) con una componente ricorrente o convoluzionale sull'asse temporale
prende il nome di \emph{spatiotemporal graph neural network} (STGNN), un
termine ormai consolidato nella letteratura sul forecasting di serie
temporali multivariate correlate \citep{cini2025graphdl}. MetaSTGAT,
descritto nel dettaglio nel Capitolo~\ref{cap:tsc-modelli}, costituisce un esempio 
di questa famiglia applicata al controllo semaforico anziché
alla previsione: un layer di \emph{graph attention} cattura la dipendenza
spaziale tra intersezioni limitrofe, un layer ricorrente (LSTM) ne
cattura l'evoluzione temporale, e i due vengono addestrati congiuntamente
end-to-end.

\paragraph{Pesi generati da meta-learning.} Sia la GCN sia la GAT, nella
loro formulazione originale, condividono gli stessi pesi su
tutti i nodi del grafo: un'assunzione ragionevole quando i nodi sono
sufficientemente omogenei, meno quando, come in una rete stradale reale, la
topologia locale di ciascun nodo (numero di strade entranti, presenza di
un'arteria, posizione periferica o centrale) può differire sensibilmente da
nodo a nodo. Un \emph{hypernetwork} affronta il problema generando i pesi
di un layer non come parametri fissi appresi direttamente, bensì come
output di una seconda rete condizionata sulle caratteristiche locali del
nodo: il Capitolo~\ref{cap:tsc-modelli} descrive nel dettaglio come
MetaSTGAT applichi questo principio sia al layer di \emph{graph attention}
sia al layer ricorrente.

\paragraph{Propagazione a lungo raggio.} Sia la GCN sia la GAT propagano
informazione a un solo passo di distanza (un arco nel grafo) per layer:
raggiungere un vicino a distanza $k$ richiede impilarne altrettanti, con un
costo di parametri crescente e un rischio noto di \emph{over-smoothing}, per il quale le
rappresentazioni dei nodi divengono troppo omogenee. Meccanismi più recenti, tra cui \textbf{SONAR} (\citealp{trenta2025sonar}), affrontano il
problema della propagazione a lungo raggio secondo un principio diverso,
ispirato a sistemi dinamici oscillatori piuttosto che alla concatenazione sequenziale di
layer: il dettaglio di questo meccanismo e la sua
applicazione al controllo semaforico sono rimandati al Capitolo~\ref{cap:tsc-modelli}.

\section{Traffic Signal Control come problema}
\label{sec:background-tsc}

Il controllo semaforico (\emph{Traffic Signal Control}, TSC) si articola
in tre approcci, in ordine crescente di complessità.

\textbf{Controllo a tempi fissi.} Una sequenza di fasi di durata
predeterminata, calcolata su stime di traffico medie, si ripete
ciclicamente senza mai osservare lo stato reale della rete. È un approccio robusto, in quanto non
può fallire per un sensore guasto, ma per costruzione non reattivo: la
stessa sequenza viene eseguita indipendentemente da congestioni anomale, 
incidenti o variazioni del traffico non modellate in fase di pianificazione.

\textbf{Controllo reattivo.} Una regola locale osserva lo stato corrente
dell'intersezione (tipicamente il numero di veicoli in coda per corsia) e
sceglie l'azione che ottimizza un criterio istantaneo. \textsc{Max-Pressure}
\citep{varaiya2013maxpressure} ne è l'esponente più studiato: modellando il
movimento dei veicoli come una rete di code \emph{store-and-forward}, definisce la pressione di una fase $\phi$ come la
differenza tra veicoli sulle corsie che quella fase servirebbe in ingresso e
in uscita, e sceglie sempre l'azione a pressione massima. Il risultato
teorico centrale dell'articolo originale è che
Max-Pressure è \emph{massimamente stabilizzante}; ciò significa che se è 
possibile stabilizzare la rete, ovvero non far divergere la lunghezza delle code nel tempo, 
Max-Pressure la stabilizza. L'algoritmo garantisce questo risultato usando solo la
lunghezza delle code sulle corsie adiacenti all'intersezione, e senza
conoscere a priori il volume di domanda in ingresso. Max-Pressure è 
tuttavia miope per costruzione: non pianifica alcun compromesso su un
orizzonte temporale non immediato, e la condizione di stabilità non
tiene in considerazione le attese tra le diverse direzioni
servite da un'intersezione, un aspetto di equità che questo lavoro esamina
esplicitamente (Capitolo~\ref{cap:confronto}).

\textbf{Controllo appreso.} L'approccio basato sul \emph{reinforcement learning} consente l'ottimizzazione 
diretta di metriche di rete globali (es. \emph{travel time} o \emph{throughput}) 
sostituendo euristiche come la pressione istantanea, al netto del 
costo computazionale dell'addestramento simulato.
Applicato a una rete di più intersezioni, il problema si formalizza come un Processo Decisionale di Markov decentralizzato a 
osservabilità parziale (Dec-POMDP): ogni intersezione osserva solo il proprio stato
locale, ma le sue azioni influenzano,
con un ritardo proporzionale alla distanza spaziale, lo stato delle intersezioni
adiacenti alterando il flusso di veicoli che le attraversano.

\paragraph{Architetture di coordinamento.} Esistono tre strategie
distinte che affrontano la dipendenza tra intersezioni vicine. Un agente
centralizzato che decide le azioni di tutta la rete con un'unica politica
elimina il problema per costruzione, ma lo spazio delle azioni congiunte
cresce esponenzialmente con il numero di intersezioni ed è impraticabile
oltre poche decine di nodi. All'estremo opposto, un agente indipendente per
intersezione, che non condivide alcuna informazione con i propri vicini;
questo approccio scala senza difficoltà ma tratta ogni intersezione come se fosse isolata,
ignorando che le sue azioni si propagano alla rete circostante. La via
intermedia, adottata anche da questo lavoro, fa comunicare agenti
indipendenti condividendo informazione locale. La difficoltà, discussa da
\citet{wei2019colight}, risiede nel fatto che la trasmissione dello stato dei vicini in 
una sequenza predefinita impone a tutte le intersezioni della rete gli stessi 
pesi interpretativi per ciascuna posizione della lista. 
Questo approccio non tiene conto della reale importanza che ogni specifico 
nodo adiacente può rivestire per le diverse intersezioni del sistema.
Una GNN evita il problema: pesando il contributo di ogni vicino in
base al suo contenuto invece che alla sua posizione in un elenco, permette a
un unico insieme di pesi condiviso di generalizzare tra intersezioni con
topologie locali diverse, come quelle della rete adottata in
questo lavoro (Capitolo~\ref{cap:tsc-modelli}).

\paragraph{Coordinamento oltre l'adiacenza fisica.} Il grafo usato 
da MetaSTGAT e da questo lavoro coincide con l'adiacenza fisica
della rete stradale: un'intersezione comunica solo con le intersezioni
immediatamente collegate da una strada. Dalla letteratura si osserva
che interazioni tra sole intersezioni adiacenti non siano sufficienti a catturare tutte 
le dinamiche di traffico e si propongono diverse soluzioni per ovviare a questo limite. 
In questo lavoro si è deciso di non affrontare questo tema, esplorandolo solo in parte nel Capitolo ~\ref{cap:confronto}.
Il problema resta comunque aperto e framework che combinano ingredienti proposti da MetaSTGAT, 
quali \emph{graph attention}, una componente ricorrente sul tempo e reinforcement
learning per il controllo semaforico, continuano ad essere oggetto di ricerca.

\paragraph{Una tassonomia delle soluzioni recenti.} La letteratura più
recente sul reinforcement learning applicato al TSC classifica i metodi
lungo assi in parte indipendenti da quelli discussi finora. Un primo asse
distingue, all'interno del multi-agente, due architetture di
coordinamento: l'\emph{addestramento centralizzato con esecuzione
decentralizzata} (CTDE), in cui gli agenti condividono informazione
globale solo durante il training per poi decidere localmente, e
l'\emph{addestramento ed esecuzione entrambi decentralizzati} (DTDE), in
cui l'informazione resta sempre e solo locale, in entrambe le fasi. 
Un secondo asse raccoglie le tecniche di apprendimento
integrate nel ciclo RL, ciascuna introdotta per rispondere a un limite
specifico \citep{xiao2025tscreview}:
reti neurali a grafo per ottenere una migliore coordinazione tra intersezioni, 
componenti ricorrenti per catturare la dipendenza tra passi successivi a livello temporale,
transfer-learning per generalizzare rapidamente a scenari nuovi, 
meta-learning per adattarsi a condizioni operative diverse,
reti generative avversarie per compensare dati di stato mancanti o incompleti.
Nessuna di queste tecniche risolve da sola il problema: la letteratura
le introduce quasi sempre in risposta a un limite preciso.


\paragraph{Dove si colloca MetaSTGAT.} Il modello di questo lavoro è la combinazione 
di alcune delle tecniche appena citate. 
La sua è un'architettura DTDE con scambio locale di informazione, che 
affianca al \emph{graph attention} (GAT) un componente ricorrente (LSTM) per catturare la dipendenza temporale oltre
a quella spaziale, collocandosi nella famiglia delle reti spatio-temporali
su grafo (STGNN); il modello sfrutta inoltre il meta-learning,
con l'obiettivo di adattare dinamicamente la propria architettura alle condizioni della rete. 

\paragraph{Simulazione.} Valutare una politica di controllo semaforico su
una rete reale non è praticabile durante lo sviluppo. Affinchè i risultati siano 
informativi, serve un simulatore che riproduca la dinamica del traffico 
con un dettaglio sufficiente, pur mantenendo un costo
computazionale accettabile per l'addestramento di centinaia di episodi. 
\textsc{CityFlow} \citep{zhang2019cityflow} è il simulatore scelto
in questo lavoro: rispetto ad alternative più diffuse come SUMO
\citep{lopez2018sumo}, privilegia
la velocità di simulazione rispetto al realismo microscopico del singolo
veicolo, un compromesso adatto quando l'obiettivo primario è addestrare un
agente su un numero elevato di episodi piuttosto che riprodurre con
precisione millimetrica la dinamica di un incrocio specifico. Il
Capitolo~\ref{cap:tsc-modelli} descrive nel dettaglio come questo lavoro
configura l'ambiente CityFlow.
